\chapter{Lab Environment, Hardware, and Checkpoints}\label{app:setup}

This appendix collects the practical details that the labs assume: how the code is laid out,
how to install it, which hardware is enough for which lab, and how to get the pretrained weights
that the measurement sections use.

\section{Repository layout}

\begin{lstlisting}[style=plain]
code/
  llmsys/             shared package, built up chapter by chapter
    layers.py         norms, MLPs, RoPE, masks, MHA, GQA
    kvcache.py        dynamic and static KV caches
    gpt.py qwen3.py   model definitions + Hugging Face weight loaders
    tokenizer.py      word-level and byte-level BPE tokenizers
    data.py train.py  datasets, loaders, training loop, LR schedule
    generate.py       sampling filters, cached / batched generation, log-probs
    flops.py          napkin-math calculators (params, FLOPs, bytes, rooflines)
    ...               later chapters add eval/, rl/, serving/, dist/, kernels/, ...
  labs/labNN_*/       README.md, starter.py, solution.py, test_labNN.py, run/bench scripts
                      (lab NN belongs to chapter NN; labA/labB to appendices A/B)
  projects/pK_*/      generic infrastructure for the capstone projects (Part VII)
  run_all_tests.sh    runs the core tests and each lab in its own process
  tests/              tests for the core package
  data/               small, original sample datasets used by the tests
\end{lstlisting}

The core modules (\texttt{layers}, \texttt{kvcache}, \texttt{gpt}, \texttt{qwen3},
\texttt{tokenizer}, \texttt{data}, \texttt{train}, \texttt{generate}) all use one model
interface:
\begin{lstlisting}
logits = model(idx, cache=None, pad_mask=None)   # idx: (B, T) -> logits: (B, T, V)
\end{lstlisting}
Passing a \code{KVCache} turns the same forward pass into an incremental, cached decode step
(\cref{ch:inference}). Passing \code{pad\_mask} (True for real tokens) masks out the
left padding used for batched generation (\cref{ch:serving}).

\section{Installation}

\begin{lstlisting}[style=shell]
python -m venv .venv && source .venv/bin/activate
pip install -r code/requirements.txt         # torch, tokenizers, safetensors, sympy, ...
cd code && python -m pytest tests            # core package self-test (CPU, ~10 s)
LAB_IMPL=solution ./run_all_tests.sh         # every lab against its reference solution
\end{lstlisting}
Run labs one directory at a time, or with \texttt{run\_all\_tests.sh}, rather than as a single
\texttt{pytest labs} invocation. The Triton labs use Triton's CPU interpreter, which must be
configured before \texttt{triton} is first imported, so they cannot share a process with labs that
import Triton normally.
For GPU work, install the PyTorch build that matches your CUDA driver
(\url{https://pytorch.org/get-started/locally/}). Triton ships with recent Linux CUDA builds of
PyTorch. The labs in \cref{app:triton} need an NVIDIA or AMD GPU, and their tests skip
automatically on CPU.

\section{Which hardware for which lab}

\begin{table}[h]
\centering\small
\begin{tabularx}{\linewidth}{@{}lX@{}}
\toprule
\textbf{Hardware} & \textbf{What runs comfortably}\\
\midrule
Laptop CPU & All tests; tiny-model training and RL demos; tokenizer, napkin-math, engine-simulation labs\\
Apple-silicon laptop (MPS) & GPT-2 small pretraining and fine-tuning; Qwen3-0.6B inference at short contexts\\
Single 16--24\,GB GPU & Everything up to GRPO and distillation on Qwen3-0.6B with short responses; LoRA on 1.7B\\
Single 80\,GB GPU (A100/H100) & Full-length GRPO runs on 0.6B--1.7B; inference on 8B; realistic serving-engine benchmarks\\
2--8 GPUs & The real-hardware variants of the distributed-training (\cref{ch:distributed}) and RL-systems (\cref{ch:rlsys}) labs\\
\bottomrule
\end{tabularx}
\caption{Rough guide to the hardware each lab needs.}
\end{table}

\begin{napkinmath}[Will it fit?]
Inference needs roughly the weights plus the KV cache: for Qwen3-0.6B in bf16 that is
about $0.6\times10^9\times 2\approx 1.2$\,GB of weights plus
$2\cdot L\cdot G\cdot \dhead\cdot 2\,\text{bytes} = 2\cdot28\cdot8\cdot128\cdot2 \approx 115$\,KB
of cache per token, so about 0.47\,GB for a 4096-token response. Full fine-tuning with AdamW
needs about 16 bytes per parameter before activations (\cref{ch:pretrain}), or about 9.6\,GB for
0.6B parameters. Use \code{llmsys.flops} to repeat these estimates for any configuration.
\end{napkinmath}

\section{Getting checkpoints and datasets}

The tests never touch the network. The \texttt{run.py}/\texttt{bench.py} scripts expect
checkpoints that you download once:

\begin{lstlisting}[style=shell]
pip install huggingface_hub
huggingface-cli download openai-community/gpt2 --local-dir ckpt/gpt2
huggingface-cli download Qwen/Qwen3-0.6B-Base --local-dir ckpt/qwen3-0.6b-base
huggingface-cli download Qwen/Qwen3-0.6B --local-dir ckpt/qwen3-0.6b   # reasoning/chat
\end{lstlisting}

\begin{lstlisting}[caption={Loading the downloaded checkpoints into our from-scratch modules.},label={lst:setup:load}]
from safetensors.torch import load_file
from llmsys.gpt import GPTModel, gpt2_config, load_gpt2_hf_state_dict
from llmsys.qwen3 import load_qwen3_from_dir, Qwen3Tokenizer

gpt2 = GPTModel(gpt2_config("gpt2-small (124M)"))
load_gpt2_hf_state_dict(gpt2, load_file("ckpt/gpt2/model.safetensors"))

qwen = load_qwen3_from_dir("ckpt/qwen3-0.6b-base", size="0.6B", device="cuda")
tok = Qwen3Tokenizer("ckpt/qwen3-0.6b-base/tokenizer.json")
\end{lstlisting}

Datasets used by the measurement scripts include MATH-500 and the MATH training split
(\cref{ch:eval,ch:grpo,ch:distill}), the SMS Spam Collection (\cref{ch:classify}), an
Alpaca-style instruction set (\cref{ch:instruct}), and any public-domain short story for
pretraining (\cref{ch:pretrain}). Each script's \texttt{--help} names the source it expects
and where to put the file. Small sample files written for this book live in \texttt{code/data/}
so the tests can run offline.

\section{Reproducibility and measurement hygiene}

\begin{itemize}
  \item Call \code{llmsys.utils.set\_seed}. On GPU, also set
        \code{torch.backends.cudnn.deterministic=True} when you compare numerics. Remember that
        batched and unbatched runs can differ in the last bits because reduction order changes.
  \item Time GPU code only after \code{torch.cuda.synchronize()} (the \code{llmsys.utils.timer}
        context manager does this). Discard warm-up iterations, which include compilation,
        autotuning and allocator growth.
  \item Report the configuration with every number: hardware, dtype, batch size, sequence
        length, PyTorch version, and whether \code{torch.compile} or CUDA graphs were used.
  \item Prefer distributions to single numbers. Report medians and percentiles (p50/p90/p99)
        for latencies, and confidence intervals for accuracies (\cref{ch:eval}).
\end{itemize}
